\section{Qiushi Engine 与自主数学研究}\label{sec:research}

\subsection{研究架构与计算环境}

Qiushi Engine 是面向持续自主科学研究的分层多智能体研究系统。其首篇关于真实光学平台上自主发现的报告，介绍了负责规划、方法发展、实验和批判性审查的核心研究层，以及知识检索、记忆、辅助探索和证据检查等支持能力\cite{yang2026optical}。经过整理的知识与记忆使早期发现能够影响后续决策，而不必反复重建全部上下文。

本研究将这些能力用于精确数学。Qiushi Engine 端到端地开展了数学研究：提出候选路线，编写并运行研究程序，比较不同表示，推导和检验命题，并形成、记录最终证明。研究最初针对更广泛的秩 22 构造与秩 23 最优性问题；随后聚焦二元域下界，将研究集中到可以精确、有限表示的问题，同时保留构造和下界两条路线。科学报告于 9 月 9 日完成。

计算环境将 Python、C++ 程序与数值线性代数、符号计算、组合搜索及证明检查结合起来。公开的研究程序包括 NumPy、SymPy 计算和约束规划模型。最终有限证明使用 PySAT 编码、CaDiCaL 生成证明、\texttt{drat-trim} 验证，以及独立的 C++ 商下界表重建程序。这些工具承担不同职责：数值实验可以提示结构；精确恒等式或经过检查的有限证明则确立结构。维护版本的数学复验可在 Linux 上运行，无须 Qiushi Engine 或语言模型访问。依赖与命令汇总于 \qpath{reproducibility/README.md}。

\subsection{研究记忆与证据}

Meta-Trace 是 Qiushi Engine 为长程自主科学研究设计的结构化研究记录机制，首篇 Qiushi Engine 论文介绍了其设计与应用\cite{yang2026optical}。它连接三个互补视角：\emph{历史}保存问题、路线与行动的摘要；\emph{证据}将命题对应到论证、计算和材料；\emph{轨迹}保留理解的变化、决策与未解决问题。分层摘要和检索机制将既往研究的概要与详细记录、支撑材料关联起来，使系统能够在需要时恢复上下文。因此，后续推理既能重新考察一个发现，也能追溯其成立条件、复用方法，或识别已经纠正过的错误。Meta-Trace 为记忆提供这一研究历史基础；底层证据、证据摘要以及由此得出的结论仍是不同的对象。

最终证明按逻辑依赖排列命题，而研究发现这些依赖的顺序通常不同。一个有希望的候选可能失败；一个较弱的表述可能揭示更强表述必须保留的信息；看似计算上的障碍，也可能转化为代数结构问题。理解这一发展过程，需要在保留结果的同时保留决策理由。

附录\ref{app:research-record}中的研究历程，是以数学问题为主线，对原始 Meta-Trace 记录及相关笔记、规划、程序和结果的综合。重复记录经过压缩，相关研究线按数学问题组织，重要纠错与未闭合分支则明确保留。研究导读将各主题连接到具体支撑材料。这种组织保留了研究中的科学联系，而不复现私有运行细节。

对 Qiushi Engine 而言，这些记录使研究保持连续，让已积累的认识能够被后续推理复用。对读者而言，它们使方法的发展可以被考察：哪项观察改变了问题，哪个反例否定了推断，哪种表示使新的证明成为可能。记录还将有用的中间方法与开放分支保存为后续研究的起点。因此，最终定理与通向定理的路径成为两个相互关联的科学对象，各自都有支撑证据。

公开这一研究历程，使贡献超出最终定理本身。数学理解的发展由此成为可供研究的对象：证据如何改变猜想，为什么放弃一个看似有前景的表示，以及中间结果如何促成另一种证明策略。研究者可以考察这些转折、复用所形成的方法，并基于同一证据比较不同路线。这样的开放记录，为 Qiushi Engine 的长程自主研究能力提供了具体、实质性的案例；能力体现在研究产生的数学决策与结果中，而不只体现在最终答案上。

三次变化构成这一历程的主线。数值延拓使精确性与系数增长成为不可分割的问题；商张量搜索揭示了合法第一因子支撑与完整分解之间的差距；最终，秩下界中的取等关系恢复了三个因子之间的联系。这是积累的认识改变研究策略的三个关键节点。下文解释这些转折；附录进一步展开中间路线及其证据。

\subsection{数值形变与精确性}

Laderman 的 23 次乘法构造，将逐元素计算的 27 个乘积重新组织为可共享的双线性计算\cite{laderman1976noncommutative}。从 23 降至 22，要求找到另一个精确恒等式。22 项分解有 $22(9+9+9)=594$ 个系数，且必须满足 729 个三次系数恒等式。因此，小矩阵维度背后隐藏着一个高度耦合的代数问题。

Qiushi Engine 探索了已知秩 23 方案的形变、配对缺陷表述，以及循环和转置对称性。某些数值残差下降时，因子幅度却持续增大：大项彼此抵消，而没有趋近有限系数解。这使问题从“残差还能降到多小”转向“何种几何结构能够支撑精确的更短分解”。关联条件表述和局部相容性检验明确了这一区别，提供了可复用的表示与检查方法，但秩 22 构造仍未解决。

\subsection{商核心与因子相容性}

在 $\FF$ 上，精确算术消除了存在性问题中的数值判零困难。商去 $E_{11}$ 方向得到的余维一核心，提供了一个具体目标：若能给出该核心的 19 项分解，将其嵌入坐标补空间，再补回被删去的三个乘积，即可得到完整张量的 22 项分解。

占据下界约束了候选的可能第一因子。然而，一个合法的第一因子多重集并不提供相容的第二、第三因子。Qiushi Engine 研究了固定因子补全、收缩秩、提升相容性、编码理论解释及矩松弛，以理解这些缺失信息。每种表示都提供了不同的相容性检验方式；若干仅表达必要条件的模型仍然可行，或在计算上未获解决。

关键认识在于：确定因子的分布，与确定各因子如何相容，是两个不同的问题。这促使研究回到完整张量：能否不先解决整个核心问题，而由支撑几何强迫出一种构型，使额外的代数恒等式不可避免？

\begin{figure}[tbp]
  \centering
  \makebox[\linewidth][c]{\begin{tikzpicture}[x=1cm,y=1cm,
  route/.style={draw=ReportRule,fill=white,align=center,text width=6.2cm,
    inner sep=8pt,font=\sffamily\small,minimum height=1.45cm},
  flow/.style={-{Stealth[length=2.2mm]},thick,draw=ReportPrimary}]
  \node[route] (broad) at (0,0)
    {\textbf{Exact algorithm search}\\Known rank-23 schemes;\\deformation and symmetry};
  \node[route] (core) at (7.15,0)
    {\textbf{Finite-field restrictions}\\Core supports; occupation;\\factor completion};
  \node[route] (lesson) at (0,-2.15)
    {\textbf{Mathematical lessons}\\Finite exactness;\\support versus factor compatibility};
  \node[route] (full) at (7.15,-2.15)
    {\textbf{Return to the full tensor}\\High-rank pairs; affine cosets;\\rank-sum equality};
  \node[route,text width=13.35cm,draw=ProofGreen] (result) at (3.575,-4.3)
    {\textbf{Structural lower-bound proof over $\FF$}\\Eight certified quotient raises
    $\longrightarrow$ affine profile $(16,1,3)$\\
    Saturation $\longrightarrow$ cross-factor identities $\longrightarrow$ contradiction};
  \draw[flow] (broad) -- (lesson);
  \draw[flow] (core) -- (full);
  \draw[flow] (lesson.east) -- (full.west);
  \draw[flow] (full.south) -- (full.south |- result.north);
\end{tikzpicture}
\unskip}
  \caption{研究中数学表示的变化。构造与商张量搜索澄清了精确性和因子相容性；回到完整张量后，这些问题与秩取等关系及结构性下界证明连接起来。}
  \label{fig:research-route}
\end{figure}

\subsection{分裂展平的饱和与代数结构}

对于假设存在的 20 项分解，分裂展平要求 $\sum_t\rk(A_t)\ge27$。经过认证的占据下界与仿射几何，强迫第一因子中恰有十六个秩一、一个秩二和三个秩三矩阵。其秩和为 $16+2+9=27$：不等式取等。

关键转折是将这一秩和等式理解为秩可加性条件。最后的证明任务因此发生变化：
\[
  P=\sum_tM_t,\qquad\sum_t\rk(M_t)=\rk(P)
  \quad\Longrightarrow\quad M_tP^{-1}M_s=\delta_{ts}M_t.
\]
取等通过 $P^{-1}$ 将第二、第三因子耦合起来。对于矩阵乘法张量，显式计算这些乘积便可证明，所必需的可逆因子无法相容。因此，第一因子计数导出了涉及全部三个因子的矛盾。最终论证用有限几何强迫取等，再由取等产生代数结构，以恒等式取代剩余的支撑枚举。

Qiushi Engine 还缩减了有限前提：实际只需要八个二维轨道，其中每个平面的三个非零矩阵都至少秩二。此前考虑的一个额外轨道，其第三方向本来就秩一。明确几何论证真正需要的命题后，删去这一前提使证明更短。

\subsection{语义验证与独立检查}

有限前提通过两种独立编码重建：以有限个布尔副本表示整数重数，以及显式的单调一元计数。其中一个商允许某个方向重数为二，因此必须区分整数计数与表示互异支撑的比特变量。此前的一次辅助变量冲突曾产生错误的矛盾。修正后的方法因此既给出从数学对象到布尔赋值的显式映射，也提供经过检查的反驳证明。

跨因子恒等式也接受了直接检验：构造实际的 $27\times27$ 置换矩阵，独立于候选指标公式计算初等矩阵乘积。记录中的检查覆盖全部 $729^2=531{,}441$ 对初等简单张量，与下文的多线性推导相互补充。过程中发现的一处约定错误，在张量恒等式层面得到纠正。

这些检查也改变了数学解释，使重数、坐标以及每个蕴含的适用范围变得明确。最终证明由有限计算基础与其后的完整书面结构推导组成。

\subsection{研究发展与可复用知识}

证明与 Meta-Trace 回答不同的科学问题。证明说明 20 项分解为何不可能；Meta-Trace 解释 Qiushi Engine 如何找到使这一矛盾可被推导的表示。早期研究提供了精确对照、商构造，以及对松弛所丢失信息的理解；后续研究利用这些区别，对完整张量提出了更有效的问题。

连贯叙述与配套材料共同展示一种方法的发展：提出了哪些想法，用什么计算检验它们，为何修订某些解释，以及保留下来的认识如何影响后续研究。自主研究过程因此与数学结果一道，成为可供科学考察和复用的对象。

\begin{figure}[tbp]
  \centering
  \makebox[\linewidth][c]{\begin{tikzpicture}[x=1cm,y=1cm,
  box/.style={draw=ReportRule,fill=white,align=center,text width=6.2cm,
    inner sep=8pt,font=\sffamily\small,minimum height=1.45cm},
  flow/.style={-{Stealth[length=2.2mm]},thick,draw=ReportPrimary}]
  \node[box] (finite) at (0,0)
    {\textbf{Certified finite premises}\\Line bounds; eight 2-plane raises;\\affine-plane bounds};
  \node[box] (split) at (7.15,0)
    {\textbf{Full-tensor split flattening}\\$\sum_t\rk(A_t)\ge\rk(P)=27$\\At least four high-rank first factors};
  \draw[dashed,ReportRule] (-3.38,-1.08) -- (10.53,-1.08);
  \node[box,text width=13.35cm] (geometry) at (3.575,-2.15)
    {\textbf{High-rank factors share an affine row or column coset}\\
     Pairwise differences have rank one;\\
     the normalized lower block has rank two};
  \node[box,text width=13.35cm] (profile) at (3.575,-4.3)
    {\textbf{Affine-plane caps force the profile $(16,1,3)$}\\
     Exactly four high-rank first factors, three invertible\\
     $\sum_t\rk(A_t)=16+2+9=27$: equality in the rank bound};
  \node[box,text width=13.35cm] (sat) at (3.575,-6.45)
    {\textbf{Equality forces cross-factor identities}\\
     $M_tP^{-1}M_s=\delta_{ts}M_t$\\
     $M_tP^{-1}M_s=M(A_tB_sC_t^TA_s,B_t,C_s)$};
  \node[box,text width=13.35cm,draw=ProofRed] (end) at (3.575,-8.6)
    {\textbf{Contradiction: at most one invertible first factor, but three are required}\\
     Diagonal identities force $B_t,C_t$ to be invertible when $A_t$ is invertible\\
     Off-diagonal identities would make an invertible product zero};
  \draw[flow] (finite.south) -- (finite.south |- geometry.north);
  \draw[flow] (split.south) -- (split.south |- geometry.north);
  \draw[flow] (geometry) -- (profile);
  \draw[flow] (profile) -- (sat);
  \draw[flow] (sat) -- (end);
\end{tikzpicture}
\unskip}
  \caption{两类互补约束在取等处汇合。有限前提一侧约束第一因子的几何，展平一侧提供秩预算。在被强迫的秩型处，预算恰好用尽，由此得到涉及全部三个因子的符号矛盾。虚线上方是有限前提与一般秩下界；下方是它们对 $\FF$ 上假设存在的 20 项分解所产生的结构性推论。}
  \label{fig:proof-structure}
\end{figure}
\FloatBarrier
